\documentclass[10pt]{extarticle}
\usepackage[a4paper,margin=1.5cm]{geometry}
\usepackage[T1]{fontenc}
\usepackage{amsmath,amssymb}
\usepackage{graphicx}
\usepackage[export]{adjustbox}
\usepackage{multicol}
\usepackage{enumitem}
\usepackage{booktabs,array,tabularx}
\setlength{\columnsep}{0.7cm}
\setlength{\parindent}{0pt}
\setlength{\parskip}{3pt}
\pagestyle{empty}
\makeatletter
% figure registry (figdefs.tex is generated by the extract step)
\newcommand{\qfigdef}[4]{\@namedef{qf@file@#1}{#2}\@namedef{qf@w@#1}{#3}\@namedef{qf@wide@#1}{#4}}
\newcommand{\qf@check}[1]{\@ifundefined{qf@file@#1}{\PackageError{score2}{Unknown figure id #1}{}}{}}
\newcommand{\qf@img}[4][]{\includegraphics[width=#2\dimexpr\@nameuse{qf@w@#4}pt\relax,max width=#3,max height=0.45\textheight,#1]{\@nameuse{qf@file@#4}}}
% Full-width questions: finalize wraps every question that contains a wide figure in \qfullbegin...\qfullend,
% so the WHOLE question (stem, figures, options) is set across the page in reading order, then the two columns
% resume. The question's head - its number, stem and first figure (\qfullsplit marks the end of the head) -
% is kept on one page, so a number is never left at a page foot with its figure overleaf; the rest may break
% across pages normally, so no big gaps are left. Without a \qfullsplit, the whole question is kept together
% if it is under 45% of a page. Never write these yourself.
\newif\ifqf@full
\newif\ifqf@boxed
\newsavebox{\qf@box}
\newcommand{\qf@open}{\global\qf@fulltrue\par\global\qf@boxedtrue
  \setbox\qf@box\vbox\bgroup\hsize\textwidth\linewidth\textwidth}
\newcommand{\qf@emit}[1]{\par\egroup\global\qf@boxedfalse
  \ifdim\dimexpr\ht\qf@box+\dp\qf@box\relax<#1\textheight\noindent\usebox{\qf@box}\par\else\unvbox\qf@box\fi}
\newcommand{\qfullsplit}{\ifqf@boxed\qf@emit{0.85}\fi}
\newcommand{\qf@close}{\ifqf@boxed\qf@emit{0.45}\fi\global\qf@fullfalse}
\newcommand{\qfullbegin}{\end{multicols}\qf@open}
\newcommand{\qfullend}{\qf@close\begin{multicols}{2}}
\newcommand{\qfullstart}{\qf@open}
\newcommand{\qfulldone}{\qf@close}
% \qfig[scale]{id}: block figure, centred, directly where it is written (size set by the tool).
% Fallback if a wide figure is not inside a full-width question: break out of the columns in place - no floats.
\newcommand{\qfig}[2][1]{\qf@check{#2}%
  \ifqf@full\qf@blk{#1}{#2}\else
  \ifnum\@nameuse{qf@wide@#2}=1 \expandafter\@firstoftwo\else\expandafter\@secondoftwo\fi
  {\end{multicols}\begin{center}\qf@img{#1}{\textwidth}{#2}\end{center}\begin{multicols}{2}}%
  {\qf@blk{#1}{#2}}\fi}
\newcommand{\qf@blk}[2]{\par\smallskip{\centering\qf@img{#1}{\linewidth}{#2}\par}\smallskip}
% \qfigi[scale]{id}: inline figure (inside option lists / tables / minipages); never breaks columns.
\newcommand{\qfigi}[2][1]{\qf@check{#2}\qf@img[valign=c]{#1}{\linewidth}{#2}}
\makeatother
% structure macros - lint reads these to check completeness and figure placement
\newcommand{\papersec}[1]{\par\bigskip{\centering\Large\bfseries PAPER-#1\par}\medskip}
\newcommand{\partsec}[2]{\par\medskip\noindent{\large\bfseries PART-#1: #2}\par\smallskip}
\newcommand{\sect}[1]{\par\smallskip\noindent\textbf{#1}\par}
\newcommand{\qn}[1]{\par\smallskip\noindent\textbf{#1.}~}
\newcommand{\paraq}[2]{\par\smallskip\noindent\textbf{Paragraph for Questions #1 and #2}\par}
\newenvironment{qoptions}{\begin{itemize}[leftmargin=1.8em,itemsep=1pt,topsep=2pt,label={}]}{\end{itemize}}
\newcommand{\opt}[2]{\item[(#1)] #2}
% figure registry (generated by score2_tools.py finalize)
\qfigdef{p1_phy_s1_q04_1}{TOAS - Test 16 (04 Sep 2022)-img/p1_phy_s1_q04_1.png}{108.2}{0}
\qfigdef{p1_phy_s2_q05_1}{TOAS - Test 16 (04 Sep 2022)-img/p1_phy_s2_q05_1.png}{108.3}{0}
\qfigdef{p1_phy_s2_q06_1}{TOAS - Test 16 (04 Sep 2022)-img/p1_phy_s2_q06_1.png}{119.0}{0}
\qfigdef{p1_phy_s2_q07_1}{TOAS - Test 16 (04 Sep 2022)-img/p1_phy_s2_q07_1.png}{118.8}{0}
\qfigdef{p1_phy_s2_q09_1}{TOAS - Test 16 (04 Sep 2022)-img/p1_phy_s2_q09_1.png}{108.3}{0}
\qfigdef{p1_phy_s2_q10_1}{TOAS - Test 16 (04 Sep 2022)-img/p1_phy_s2_q10_1.png}{134.7}{0}
\qfigdef{p1_phy_s3_para1_2_1}{TOAS - Test 16 (04 Sep 2022)-img/p1_phy_s3_para1_2_1.png}{109.4}{0}
\qfigdef{p1_phy_s4_q08_1}{TOAS - Test 16 (04 Sep 2022)-img/p1_phy_s4_q08_1.png}{108.2}{0}
\qfigdef{p1_phy_s4_q09_1}{TOAS - Test 16 (04 Sep 2022)-img/p1_phy_s4_q09_1.png}{156.0}{0}
\qfigdef{p1_chem_s1_q04_1}{TOAS - Test 16 (04 Sep 2022)-img/p1_chem_s1_q04_1.png}{204.4}{0}
\qfigdef{p1_chem_s2_q08_1}{TOAS - Test 16 (04 Sep 2022)-img/p1_chem_s2_q08_1.png}{209.4}{0}
\qfigdef{p1_chem_s2_q08_2}{TOAS - Test 16 (04 Sep 2022)-img/p1_chem_s2_q08_2.png}{50.0}{0}
\qfigdef{p1_chem_s2_q08_3}{TOAS - Test 16 (04 Sep 2022)-img/p1_chem_s2_q08_3.png}{55.8}{0}
\qfigdef{p1_chem_s2_q08_4}{TOAS - Test 16 (04 Sep 2022)-img/p1_chem_s2_q08_4.png}{60.7}{0}
\qfigdef{p1_chem_s2_q08_5}{TOAS - Test 16 (04 Sep 2022)-img/p1_chem_s2_q08_5.png}{60.7}{0}
\qfigdef{p1_chem_s2_q09_1}{TOAS - Test 16 (04 Sep 2022)-img/p1_chem_s2_q09_1.png}{151.3}{0}
\qfigdef{p1_chem_s3_q05_1}{TOAS - Test 16 (04 Sep 2022)-img/p1_chem_s3_q05_1.png}{213.0}{0}
\qfigdef{p1_chem_s3_q06_1}{TOAS - Test 16 (04 Sep 2022)-img/p1_chem_s3_q06_1.png}{60.4}{0}
\qfigdef{p1_chem_s3_q06_2}{TOAS - Test 16 (04 Sep 2022)-img/p1_chem_s3_q06_2.png}{51.9}{0}
\qfigdef{p1_chem_s3_q06_4}{TOAS - Test 16 (04 Sep 2022)-img/p1_chem_s3_q06_4.png}{52.0}{0}
\qfigdef{p1_chem_s3_q06_3}{TOAS - Test 16 (04 Sep 2022)-img/p1_chem_s3_q06_3.png}{26.1}{0}
\qfigdef{p1_chem_s3_q06_6}{TOAS - Test 16 (04 Sep 2022)-img/p1_chem_s3_q06_6.png}{59.9}{0}
\qfigdef{p1_chem_s3_q06_5}{TOAS - Test 16 (04 Sep 2022)-img/p1_chem_s3_q06_5.png}{56.5}{0}
\qfigdef{p1_chem_s3_q06_8}{TOAS - Test 16 (04 Sep 2022)-img/p1_chem_s3_q06_8.png}{116.3}{0}
\qfigdef{p1_chem_s3_q06_7}{TOAS - Test 16 (04 Sep 2022)-img/p1_chem_s3_q06_7.png}{54.6}{0}
\qfigdef{p1_chem_s3_q06_10}{TOAS - Test 16 (04 Sep 2022)-img/p1_chem_s3_q06_10.png}{40.2}{0}
\qfigdef{p1_chem_s3_q06_9}{TOAS - Test 16 (04 Sep 2022)-img/p1_chem_s3_q06_9.png}{55.1}{0}
\qfigdef{p2_phy_s1_q01_1}{TOAS - Test 16 (04 Sep 2022)-img/p2_phy_s1_q01_1.png}{160.9}{0}
\qfigdef{p2_phy_s1_q02_1}{TOAS - Test 16 (04 Sep 2022)-img/p2_phy_s1_q02_1.png}{134.7}{0}
\qfigdef{p2_phy_s1_q03_1}{TOAS - Test 16 (04 Sep 2022)-img/p2_phy_s1_q03_1.png}{129.6}{0}
\qfigdef{p2_phy_s1_q05_1}{TOAS - Test 16 (04 Sep 2022)-img/p2_phy_s1_q05_1.png}{108.3}{0}
\qfigdef{p2_phy_s2_para7_8_1}{TOAS - Test 16 (04 Sep 2022)-img/p2_phy_s2_para7_8_1.png}{134.7}{0}
\qfigdef{p2_phy_s2_para9_10_1}{TOAS - Test 16 (04 Sep 2022)-img/p2_phy_s2_para9_10_1.png}{82.1}{0}
\qfigdef{p2_phy_s3_para1_2_1}{TOAS - Test 16 (04 Sep 2022)-img/p2_phy_s3_para1_2_1.png}{203.5}{0}
\qfigdef{p2_phy_s3_para3_4_1}{TOAS - Test 16 (04 Sep 2022)-img/p2_phy_s3_para3_4_1.png}{134.9}{0}
\qfigdef{p2_phy_s3_para5_6_1}{TOAS - Test 16 (04 Sep 2022)-img/p2_phy_s3_para5_6_1.png}{114.6}{0}
\qfigdef{p2_phy_s4_q07_1}{TOAS - Test 16 (04 Sep 2022)-img/p2_phy_s4_q07_1.png}{134.7}{0}
\qfigdef{p2_phy_s4_q08_1}{TOAS - Test 16 (04 Sep 2022)-img/p2_phy_s4_q08_1.png}{119.0}{0}
\qfigdef{p2_chem_s1_q06_1}{TOAS - Test 16 (04 Sep 2022)-img/p2_chem_s1_q06_1.png}{236.8}{0}
\qfigdef{p2_chem_s1_q06_2}{TOAS - Test 16 (04 Sep 2022)-img/p2_chem_s1_q06_2.png}{59.9}{0}
\qfigdef{p2_chem_s1_q06_3}{TOAS - Test 16 (04 Sep 2022)-img/p2_chem_s1_q06_3.png}{59.9}{0}
\qfigdef{p2_chem_s1_q06_4}{TOAS - Test 16 (04 Sep 2022)-img/p2_chem_s1_q06_4.png}{117.0}{0}
\qfigdef{p2_chem_s1_q06_5}{TOAS - Test 16 (04 Sep 2022)-img/p2_chem_s1_q06_5.png}{79.9}{0}
\qfigdef{p2_chem_s2_q08_1}{TOAS - Test 16 (04 Sep 2022)-img/p2_chem_s2_q08_1.png}{160.7}{0}
\qfigdef{p2_chem_s3_para5_6_1}{TOAS - Test 16 (04 Sep 2022)-img/p2_chem_s3_para5_6_1.png}{240.4}{0}
\raggedcolumns
\makeatletter\@beginparpenalty=10000\makeatother
\newcolumntype{L}[1]{>{\raggedright\arraybackslash}p{#1}}
\usepackage{needspace}
\let\oldsect\sect\let\oldparaq\paraq
\renewcommand{\sect}[1]{\Needspace{7\baselineskip}\oldsect{#1}\nopagebreak[4]}
\renewcommand{\paraq}[2]{\Needspace{6\baselineskip}\oldparaq{#1}{#2}\nopagebreak[4]}
\begin{document}
\begin{center}{\Large\bfseries TOAS --- Test 16 (04 Sep 2022)}\end{center}
\begin{multicols}{2}
\raggedcolumns
\papersec{1}
\partsec{1}{PHYSICS}
\sect{SECTION-I(i)}
\qn{1} At $t=0$, an electron $(e,m)$ leaves one plate of a parallel plate capacitor with negligible initial velocity. An accelerating voltage, varying as $V=kt$, is applied between the plates. The separation between the plates is $\ell$, velocity with which the electron reaches the opposite plate is :
\begin{qoptions}
\opt{A}{$\dfrac{ek}{m\ell}\left[\dfrac{3m\ell^2}{ek}\right]^{2/3}$}
\opt{B}{$\dfrac12\dfrac{ek}{m\ell}\left[\dfrac{6m\ell^2}{ek}\right]^{2/3}$}
\opt{C}{$\dfrac{2ek}{m\ell}\left[\dfrac{6m\ell^2}{ek}\right]^{2/3}$}
\opt{D}{$\dfrac{ek}{m\ell}\left[\dfrac{6m\ell^2}{ek}\right]^{1/3}$}
\end{qoptions}
\qn{2} The potential of a certain electric field in the 3-D space is given as $V=\alpha(xy-z^2)$, $\alpha$ is a constant. The projection of $\vec E$ on to the direction of vector $\vec A=\hat i+\hat j+3\hat k$ at a point $(2,1-3)$ is:
\begin{qoptions}
\opt{A}{Zero}
\opt{B}{$\dfrac{-21\alpha}{\sqrt{11}}$}
\opt{C}{$-\sqrt{11}\,\alpha$}
\opt{D}{$\dfrac{1}{\sqrt{11}}$}
\end{qoptions}
\qn{3} The function $y(x,t)=(15.0\text{ cm})\cos(\pi x-15\pi t)$, with $x$ in meters and $t$ in seconds, describes a wave on a taut string. What is the transverse speed for a point on the string at an instant when that point has the displacement $y=+12.0$ cm ?
\begin{qoptions}
\opt{A}{$180\pi$ cm/s}
\opt{B}{$225\pi$ cm/s}
\opt{C}{$45\pi$ cm/s}
\opt{D}{$135\pi$ cm/s}
\end{qoptions}
\qn{4} A current $I$ flows in a long single layer solenoid with cross-sectional radius $R$. Number of turns per unit length of the solenoid is equal to $n$. Find the maximum possible value of current $I$ that can be maintained without rupturing, if $F_0$ is the maximum tensile force the wire of solenoid can withstand.
\qfig{p1_phy_s1_q04_1}
\begin{qoptions}
\opt{A}{$I_{\max}=\sqrt{\dfrac{F_0}{\mu_0nR}}$}
\opt{B}{$I_{\max}=\sqrt{\dfrac{F_0}{2\mu_0nR}}$}
\opt{C}{$I_{\max}=\dfrac12\sqrt{\dfrac{F_0}{\mu_0nR}}$}
\opt{D}{$I_{\max}=\sqrt{\dfrac{2F_0}{\mu_0nR}}$}
\end{qoptions}
\sect{SECTION-I(ii)}
\qn{5} A charged particle flies through the electric field of a cylindrical capacitor and gets into a uniform transverse magnetic field with induction $B$ as shown in figure. In the capacitor the particle moves along the arc of a circle, in the magnetic field, along a semicircle of radius `$r$'. The potential difference applied to the capacitor is equal to $V_0$, the radii of electrodes are equal to `$a$' and `$b$' with $a<b$
\qfig{p1_phy_s2_q05_1}
\begin{qoptions}
\opt{A}{The velocity of the particle is $v=\dfrac{V_0}{2Br\ell n\left(\frac ba\right)}$}
\opt{B}{The velocity of the particle is $v=\dfrac{V_0}{Br\ell n\left(\frac ba\right)}$}
\opt{C}{The specific charge ratio $\dfrac qm=\dfrac{V_0}{B^2r^2\ell n\left(\frac ba\right)}$}
\opt{D}{The specific charge ratio $\dfrac qm=\dfrac{V_0}{2B^2r^2\ell n\left(\frac ba\right)}$}
\end{qoptions}
\qn{6} For an AC circuit shown in figure, the supply voltage has constant rms value $V$. Circuit is in resonance at frequency, $f_0$. At resonance
\qfig{p1_phy_s2_q06_1}
\begin{qoptions}
\opt{A}{Current in circuit is $\dfrac VR$}
\opt{B}{$f_0=500$ Hz.}
\opt{C}{Voltage across the capacitor is $180^\circ$ out of phase with the voltage across inductor.}
\opt{D}{Current in circuit is $\dfrac V{2R}$.}
\end{qoptions}
\qn{7} Figure shows conducting and smooth fixed rails in horizontal plane. An inductor is connected between rails. There exists a uniform constant vertical magnetic field $B$. Slider having mass $m$ length $\ell$ was initially at rest. Now a constant force $F$ is applied on slider as shown in figure. Then :
\qfig{p1_phy_s2_q07_1}
\begin{qoptions}
\opt{A}{Angular frequency of slider would be $\dfrac{B\ell}{\sqrt{mL}}$.}
\opt{B}{Angular frequency of slider would be $\dfrac{B\ell}{mL}$.}
\opt{C}{Slider performs SHM}
\opt{D}{Slider moves with constant acceleration.}
\end{qoptions}
\qn{8} A rope, under a tension of 200 N and fixed at both ends, oscillates in a second-harmonic standing wave pattern. The displacement of the rope is given by $y=(0.10\text{m})(\sin\pi x/2)\sin12\pi t$, where $x=0$ at one end of the rope, $x$ is in meters, and $t$ is in seconds.
\begin{qoptions}
\opt{A}{The length of the rope is 6m.}
\opt{B}{Mass/length of the rope is $\dfrac{25}{72}$ kg/m .}
\opt{C}{There is an antinode at $x=3$m.}
\opt{D}{The particle at $x=1$m \& $x=1.5$ m always oscillate in same phase.}
\end{qoptions}
\qn{9} A block of mass $3m$ is suspended vertically by a rope of mass $m$ and length $\ell$ as shown in figure. A wave pulse of wavelength $\lambda_0$ is generated at bottom most point A of the string.
\qfig{p1_phy_s2_q09_1}
\begin{qoptions}
\opt{A}{Wavelength of wave pulse at point B is $\sqrt{\dfrac34}\,\lambda_0$.}
\opt{B}{Wavelength of wave pulse at point B is $\sqrt{\dfrac43}\,\lambda_0$.}
\opt{C}{Time taken by wave pulse to reach point B is $\left(2-\sqrt3\right)\sqrt{\dfrac\ell g}$}
\opt{D}{Time taken by wave pulse to reach point B is $2\left(2-\sqrt3\right)\sqrt{\dfrac\ell g}$.}
\end{qoptions}
\qn{10} System shown at rest is slightly compressed and released at $t=0$. Separation between the blocks is maximum for first time at $t=t_1$ and relative velocity between the blocks is maximum at $t=t_2$, then
\qfig{p1_phy_s2_q10_1}
\begin{qoptions}
\opt{A}{$t_1=1$s}
\opt{B}{$t_1=\dfrac12$s}
\opt{C}{$t_2=1$s}
\opt{D}{$t_2=\dfrac12$s}
\end{qoptions}
\sect{SECTION-II(i)}
\paraq{1}{2}
A safety engineer needs to evaluate the practice of storing inflammable conducting liquids in non-conducting containers. The company supplying the liquid uses a cylindrical plastic container of radius $r=20$ cm and fills it to a height of $h=20$ cm, which is not equal to its full height. During handling, the exterior surface acquires a charge density of $\sigma=2\times10^{-6}$ C/m$^2$, approximately uniform. This induces a charge separation within the liquid, because it is conducting (see figure).
\qfig{p1_phy_s3_para1_2_1}
\qn{1} How much charge is induced within the bulk of the liquid (in $\mu$C) Take : $\pi\approx3$.
\qn{2} Assume that the capacitance of the central part relative to ground is nearly 36pF. Calculate the potential energy stored in the capacitor (in mJ).
\paraq{3}{4}
An accumulator of emf 2V and negligible internal resistance is connected across a uniform wire of length 10 m and resistance $30\Omega$. The appropriate terminals of a cell of emf 1.5 V and internal resistance $1\Omega$ is connected to one end of the wire and the other terminal of the cell connected through a sensitive galvanometer to a slider on the wire.
\qn{3} What is the length (in m) of the wire that will be required to produce zero deflection of the galvanometer.
\qn{4} What will be the balancing length (in m) when a coil of resistance $5\,\Omega$ is placed in series with the accumulator.
\paraq{5}{6}
A circuit takes a current of 3 A at a power factor of 0.6 lagging when connected to a 115 V- 50Hz supply. Another circuit takes a current of 5 A at a power factor of 0.707 leading when connected to the same supply. If the two circuits are connected in series across a 230 V, 50 Hz supply, then
\qn{5} What would be the value of current (in A) drawn by the circuit in ampere.
\qn{6} The power consumed by the circuit in kw.
\sect{SECTION-II(ii)}
\qn{7} A point sound source has an output of $4\pi\mu$W. What is loudness level in (dB) at a distance of 10 m from the source ?
\qn{8} The magnetic force on current carrying wire AB due to an infinitely long current wire at O, is $\dfrac{\mu_0I^2}{2\pi}\ell n(N)$ , find N :
\qfig{p1_phy_s4_q08_1}
\qn{9} A superconducting ring of mass 0.1 kg, inductance 0.02 H and radius 0.5 m is projected origin towards a magnet with velocity $v_0=1$m/s as shown. An imaginary magnetic field along axis O varies as $B_X=B_0(1+x)$ and its radial component as $B_r=2B_0$, where $B_0=0.1$ T and $x$ is in meters. Maximum distance traversed by ring (in cm) towards magnet is $N$ cm, find $N$. (Take $\pi^2=10$)
\qfig{p1_phy_s4_q09_1}
\partsec{2}{CHEMISTRY}
\sect{SECTION-I(i)}
\qn{1} Which of the following statement is incorrect?
\begin{qoptions}
\opt{A}{In chromyl chloride test, formation of chromyl chloride is a non redox change}
\opt{B}{The brown ring test of NO$_2^-$ can be done in presence of NO$_3^-$ but brown ring test of NO$_3^-$ cannot be done in presence of NO$_2^-$}
\opt{C}{CrO$_4^{2-}$(aq.) and Cr$_2$O$_7^{2-}$(aq.) can be inter changed by modifying the temperature}
\opt{D}{Borate can be detected by using CaF$_2$ along with conc. H$_2$SO$_4$ without touching the flame}
\end{qoptions}
\qn{2} Which of the following observations do not confirm presence of NH$_3$?
\begin{qoptions}
\opt{A}{Nessler's reagent gives brown ppt}
\opt{B}{Neutralising NaOH solution}
\opt{C}{Turning green solution of NiCl$_2$ to deep blue}
\opt{D}{Turning CuSO$_4$ to deep blue}
\end{qoptions}
\qn{3} Which of the following reaction is non-redox?
\begin{qoptions}
\opt{A}{Cl$^-$ + conc. H$_2$SO$_4\to$}
\opt{B}{Br$^-$ + conc. H$_2$SO$_4\to$}
\opt{C}{Fe$^{3+}$ + S$_2$O$_3^{2-}$ $\xrightarrow[\text{standing}]{\text{on long}}$}
\opt{D}{Cu$^{2+}$ + S$_2$O$_3^{2-}\to$}
\end{qoptions}
\qn{4} What reagents (conditions) are needed for the following reaction ?
\qfig{p1_chem_s1_q04_1}
\begin{qoptions}
\opt{A}{reagent A : CH$_3$COCl/AlCl$_3$; reagent B : Sn, HCl, heat}
\opt{B}{reagent A : CH$_3$CH$_2$Cl/AlCl$_3$; reagent B : KMnO$_4$, heat}
\opt{C}{reagent A : CH$_3$COCl/AlCl$_3$; reagent B : H$_2$NNH$_2$, KOH, H$_2$O, heat}
\opt{D}{reagent A : HNO$_3$, H$_2$SO$_4$; reagent B : Sn, HCl, heat}
\end{qoptions}
\sect{SECTION-I(ii)}
\qn{5} The following process has been used to obtain iodine from oil-field brines in California :\par
\noindent NaI + AgNO$_3$ $\xrightarrow{\text{Yield}=100\%}$ AgI + NaNO$_3$\par
\noindent AgI + Fe $\xrightarrow{\text{Yield}=80\%}$ FeI$_2$ + Ag\par
\noindent FeI$_2$ + Cl$_2$ $\xrightarrow{\text{Yield}=100\%}$ FeCl$_3$ + I$_2$\par
$\dfrac{1}{2\times10^4}$ part of the sample of I$_2$ produced is dissolved in water to make 1 litre solution. 10 ml of this solution requires 20 ml 0.1N hypo solution to decolorise the blue starch iodine complex. Select the correct statement(s) :
\begin{qoptions}
\opt{A}{The mass of AgNO$_3$ initially taken is 680 g}
\opt{B}{The mass of AgNO$_3$ initially taken is 850 kg}
\opt{C}{Fe required is 2000 mole}
\opt{D}{Fe required is 2500 mole}
\end{qoptions}
\qn{6} Which of the following statement(s) is/are true for ideal gas ?
\begin{qoptions}
\opt{A}{$(-w_{\text{rev}})-(-w_{\text{irr}})>0$ for isothermal expansion.}
\opt{B}{$w_{\text{rev}}>w_{\text{irr}}$ for isothermal compression}
\opt{C}{Heat rejected in isothermal process may be greater in $(n-1)^{\text{th}}$ stage as compared to $n^{\text{th}}$ stage work.}
\opt{D}{For expansion and compression both work done by the system is maximum for reversible process.}
\end{qoptions}
\qn{7} Reaction with acidified permagnte solution of (CH$_3$)$_2$C=CH--CH$_2$--COOH gives (A) and (B) compounds. A gives postive Iodoform test\par
Choose correct staement(s) of compound (A) and (B)
\begin{qoptions}
\opt{A}{Compound (A) give 2$^\circ$--alcohol by LiAlH$_4$ in T.H.F.}
\opt{B}{Compound (B) is heated to give acetic acid}
\opt{C}{Compound (A) give Iodoform test}
\opt{D}{Compound (B) is stronger acid than oxalic acid}
\end{qoptions}
\qn{8} Predict product (A), (B) (C) and (D)?
\qfig{p1_chem_s2_q08_1}
\begin{qoptions}
\opt{A}{A is \qfigi{p1_chem_s2_q08_2}}
\opt{B}{D is \qfigi{p1_chem_s2_q08_3}}
\opt{C}{B is \qfigi{p1_chem_s2_q08_4}}
\opt{D}{A is \qfigi{p1_chem_s2_q08_5}}
\end{qoptions}
\qn{9} \qfigi{p1_chem_s2_q09_1}\par
Choose incorrect statement for compound (X) and (Y)
\begin{qoptions}
\opt{A}{Compound (X) is chiral}
\opt{B}{Compound (Y) is 2--oxo--cyclopentanecarboxylic acid}
\opt{C}{Compound (Y) is react with hydroxyl amine to give two geometrical isomers}
\opt{D}{Degree of unsaturation of (X) is 3}
\end{qoptions}
\qn{10} Which of the following statement is/are correct towards electrophilic substitution reaction when given groups are attached to benzene ring-
\begin{qoptions}
\opt{A}{--NO$_2$ is a metadirecting and deactivating group}
\opt{B}{--Cl is a metadirecting and deactivating group}
\opt{C}{--OMe is an ortho-paradirecting and activating group}
\opt{D}{--CH$_3$ is an ortho-paradirecting and activating group}
\end{qoptions}
\sect{SECTION-II(i)}
\paraq{1}{2}
In qualitative analysis, the observations of each test must be carefully observed along with pH conditions and temperature conditions of the solution.
\qn{1} How many of the following steps/observations is/are wrong regarding brown ring test of NO$_3^-$ ?\par
(i) Ferrous sulphate can be allowed to stand before mixing it with the solution\par
(ii) A dark brown ring of nitroso ferrous solution is formed at the junction of the solution\par
(iii) In 1 ml of aqueous solution of salt, 2 mL of dil H$_2$SO$_4$ needs to be added\par
(iv) In case of presence of NO$_2^-$ along with NO$_3^-$, use of sulphamic acid can help with the test.\par
(v) The complex formed is coloured due to charge transfer spectrum\par
(vi) In case of formation of white ppt. while mixing conc. H$_2$SO$_4$, the test can be carried out without any measures.
\qn{2} For the complex formed in sodium nitroprusside test of sulphide ion,\par
\par\noindent\hspace*{0.5em}$=\dfrac{\parbox{0.8\linewidth}{\centering\footnotesize Difference between magnitude of charge on product complex ion and reactant complex ion}}{\parbox{0.8\linewidth}{\centering\footnotesize Number of linkage isomers if ambidentate behaviour of CN$^-$ is not considered}}$
\paraq{3}{4}
Consider the following elementary reaction X$_2$(g) $\underset{K_b=4\,\text{sec}^{-1}\text{torr}^{-1}}{\overset{K_f=4\,\text{sec}^{-1}}{\rightleftharpoons}}$ 2X(g) at $2000^\circ$C and at 2 torr equilibrium pressure. The composition of this equilibrium mixture is determined by measuring the rate of effusion through a pin hole. It is observed that the mixture effuses 1.5 times slower than helium under identical conditions.
\qn{3} Calculate molar mass of the gas X$_2$ in gram per mol.
\qn{4} If above elementary reaction is assumed to be irreversible then calculate time in sec in which amount of reactant reacted is exactly equal to the amount of reactant reacted to reach equilibrium when reaction was occuring reversibly. $\left(\text{Given : }\ell n\dfrac32=0.4\right)$
\paraq{5}{6}
Activating groups in benzene are more reactive and deactivating groups in benzene are less reactive.
\qn{5} How many of the following compounds will show higher rate of electrophilic substitution reaction than benzene -
\qfig{p1_chem_s3_q05_1}
\qn{6} How many reactions giving correct product?\par
\begin{center}\footnotesize\setlength{\tabcolsep}{2pt}
\begin{tabular}{@{}c@{\,}c@{}}
\qfigi{p1_chem_s3_q06_1} & \qfigi{p1_chem_s3_q06_2}\\
\qfigi{p1_chem_s3_q06_4} & \qfigi{p1_chem_s3_q06_3}\\
\qfigi{p1_chem_s3_q06_6} & \qfigi{p1_chem_s3_q06_5}\\
\qfigi{p1_chem_s3_q06_8} & \qfigi{p1_chem_s3_q06_7}\\
\qfigi{p1_chem_s3_q06_10} & \qfigi{p1_chem_s3_q06_9}\\
\end{tabular}\end{center}
\sect{SECTION-II(ii)}
\qn{7} How many of the following processes involves blue colouration?\par
(1) Cr$_2$O$_7^{2-}$ + H$_2$O$_2$ $\xrightarrow{\text{amyl alcohol}}$\par
(2) Fe$^{2+}$ + [Fe(CN)$_6$]$^{4-}\to$\par
(3) NH$_3$(excess) + NiCl$_2$(aq.)$\to$\par
(4) CuCl$_2$ dil solution in water\par
(5) Green solution of K$_2$MnO$_4$ is presence of acid\par
(6) Cu + ZnSO$_4$ (aq.)$\to$\par
(7) [Co(H$_2$O)$_6$]$^{2+}$ + Cl$^-\to$\par
(8) Fe$^{3+}$ + [Fe(CN)$_6$]$^{4-}\to$\par
(9) NO$_2$ + NO$_2$ $\xrightarrow{\text{excess cooling}}$\par
(10) Al$_2$(SO$_4$)$_3$ + Na$_2$CO$_3$ + Co(NO$_3$)$_2$ $\xrightarrow[\text{fusion in absence of O}_2]{\Delta}$
\qn{8} Let `A' and `B' are two completely miscible liquids. $P_A^\circ=90$ torr and $P_B^\circ=40$ torr. In a flask 8 moles of `A' is mixed with 4 moles of `B'. However as soon as B is added `A' starts dissociating as two different liquids `C' and `D' as A $\to$ C + D.\par
Both `C' and `D' are also completely miscible with each other and with `A' and `B'. Given $P_C^\circ=50$ torr and $P_D^\circ=40$ torr.\par
Assume dissociation of `A' follows 1$^{\text{st}}$ order kinetics with half life equal to 50 min. After 100 min 2 moles of a non-volatile non-electrolyte solute `S' is added in solution which immediately arrests the dissociation of A completely.\par
Calculate final vapour pressure of the solution in torr.
\qn{9} Alpha particles capture free moving electrons having wavelength 5\AA\ and form excited H-like species. In this excited state, the circumference of orbit is 2 times the wavelength of electron.\par
If photons emitted in this process fall on a metal surface of work function 10.6 eV. Find the maximum possible kinetic energy (in eV) of photoelectrons.
\partsec{3}{MATHEMATICS}
\sect{SECTION-I(i)}
\qn{1} Let $f:[1,\infty)\to(-\infty,-1]$ be a continuous and differentiable function such that $f(x)=x-2x^2+\displaystyle\int_1^x\frac{f(t)}{x}\,dt$ for all $x\in[1,\infty)$. Then, which of the following statement(s) is(are) True ?
\begin{qoptions}
\opt{A}{The curve $y=f(x)$ passes through the point $(3,-24)$}
\opt{B}{Function $f(x)$ is an onto function}
\opt{C}{Function $f(x)$ is an into function}
\opt{D}{The area of the region $\{(x,y)\in[1,2]\times R:f(x)<y<-x\}$ is $\dfrac72$}
\end{qoptions}
\qn{2} Let $A$ be a $3\times3$ matrix. If every column vector $X$ satisfies $X^TAX=0$ and $a_{12}=2008$, $a_{13}=1010$, $a_{23}=-2012$. Then the value of $a_{21}+a_{31}+a_{32}$ is (where $X^T$ represents transpose of matrix $X$)
\begin{qoptions}
\opt{A}{$-6$}
\opt{B}{1006}
\opt{C}{$-1006$}
\opt{D}{0}
\end{qoptions}
\qn{3} The solution of the differential equation $\dfrac{dy}{dx}+x(x+y)=x^3(x+y)^3-1$ is : (Where $c$ is arbitrary constant)
\begin{qoptions}
\opt{A}{$\dfrac{1}{x+y}=x^2+1+ce^x$}
\opt{B}{$\dfrac{1}{(x+y)^2}=x^2+1+ce^{x^2}$}
\opt{C}{$\dfrac{1}{(x+y)^2}=x+1+ce^x$}
\opt{D}{$\dfrac{1}{x+y}=x+1+ce^{x^2}$}
\end{qoptions}
\qn{4} Consider a real valued continuous function $f$ such that $f(x)=\sin x+\displaystyle\int_{-\pi/2}^{\pi/2}(\sin x+tf(t))\,dx$. If $M$ and $m$ are maximum and minimum value of the function $f$, then
\begin{qoptions}
\opt{A}{$\dfrac Mm=4$}
\opt{B}{$M-m=2\pi+1$}
\opt{C}{$M+m=4(\pi+1)$}
\opt{D}{$Mm=2(\pi^2+1)$}
\end{qoptions}
\sect{SECTION-I(ii)}
\qn{5} If $A$ and $B$ are two matrices such that $AB=A$ and $BA=B$, then
\begin{qoptions}
\opt{A}{$(A^3-B^3)^5=(A^5-B^5)^3$}
\opt{B}{$(A^4-B^4)^3=A^{12}-B^{12}$}
\opt{C}{$(A-B)$ is idempotent}
\opt{D}{$(A^{10}-B^{10})$ is nilpotent}
\end{qoptions}
\qn{6} Let $\displaystyle\int\frac{dx}{x^2+ax+1}=f(g(x))+C$, where $C$ is constant of integration, then
\begin{qoptions}
\opt{A}{$f(x)$ is inverse trigonometric function for $|a|<2$}
\opt{B}{$f(x)$ is logarithmic function for $|a|>2$}
\opt{C}{$g(x)$ is quadratic function for $|a|<2$}
\opt{D}{$g(x)$ is rational function for $|a|>2$}
\end{qoptions}
\qn{7} Let $A$ is a non-singular matrix of order $n\times n$ such that $3ABA^{-1}+A=2A^{-1}BA$, then which is/are correct (where $|A|=\det(A)$) -
\begin{qoptions}
\opt{A}{$A$ and $B$ both are identity matrix}
\opt{B}{$|A+B|=0$}
\opt{C}{$|ABA^{-1}-A^{-1}BA|=0$}
\opt{D}{$A+B$ is not a singular matrix.}
\end{qoptions}
\qn{8} Let $f(x)$ be a non-constant, twice differentiable function defined on $R$ such that $y=f(x)$ is symmetric about line $x=1$ and $f(-1)=f'\left(\dfrac14\right)=f'\left(\dfrac12\right)=0$, then which of following is/are correct statement(s) ?
\begin{qoptions}
\opt{A}{$f''(x)=0$ has atleast four roots in the interval $(0,2)$}
\opt{B}{$\displaystyle\int_{-2}^2(x^5+x^3)f(1+x)\,dx<\int_0^2\frac{1}{1+2^{f(x)}}\,dx$}
\opt{C}{there exist at least one $C\in\left(\dfrac32,\dfrac74\right)$ such that $f'(C)+Cf''(C)=0$}
\opt{D}{for some $C\in\left(1,\dfrac32\right)$, $f'(C)=Cf''(C)$}
\end{qoptions}
\qn{9} Let $f$ be twice differentiable real valued function satisfying $f'(x)\neq0$, $f(x)+f''(x)=-xg(x)f'(x)$ where $g(x)>0\ \forall\ x>0$. If $f(0)=-3$ and $f'(0)=4$ then
\begin{qoptions}
\opt{A}{$|f(x)|\le5\ \forall\ x>0$}
\opt{B}{$(f(x))^2+(f'(x))^2$ is decreasing $\forall\ x>0$.}
\opt{C}{$|f(x)|>5\ \forall\ x>0$}
\opt{D}{$(f(x))^2+(f'(x))^2$ is increasing $\forall\ x>0$.}
\end{qoptions}
\qn{10} Given $f(x)-2\displaystyle\int_0^{\pi/4}\left(\frac{\sin^2x}{\cos^5x}\right)f(t)\cos t\,dt=\frac{\sin^2x}{\cos^5x}$.
\begin{qoptions}
\opt{A}{If $f:R-\left\{(2n+1)\dfrac\pi2\right\}\to R$, then $f$ is one-one ; $n\in Z$}
\opt{B}{$f(x)$ is periodic function with fundamental period $\pi$}
\opt{C}{If $f:R-\left\{(2n+1)\dfrac\pi2\right\}\to R$, then $f$ is onto ; $n\in Z$}
\opt{D}{$f\left(\dfrac\pi3\right)=72$}
\end{qoptions}
\sect{SECTION-II(i)}
\paraq{1}{2}
Let $A=[a_{ij}]_{3\times3}$ be upper triangular matrix of order 3 and its characteristic equation is $A^3+\alpha A^2+\beta A-I=0$, where $a_{ij}\in\{-1,0,1\}$. (where $\alpha,\beta$ are scalars)
\qn{1} Sum of all the possible values of $\alpha$ is -
\qn{2} If $\alpha=-3$, then value of $\beta$ is -
\paraq{3}{4}
Let $f:[0,\infty)\to[2,\infty)$ be a derivable function which is also surjective such that $f(0)=2$ and satisfying $f^2(x)-f^2(y)=3f(x)-3f(y)+\sqrt x-\sqrt y$
\qn{3} The value of $\displaystyle\lim_{x\to0^+}\left(\frac{2f(x)-3-e^{\sqrt x}}{\sqrt x}\right)$ equals
\qn{4} If $g$ is the inverse function of $f$ then $\dfrac{d}{dx}\left(\dfrac{4}{g(x)}\right)$ at $x=3$ is equal to
\paraq{5}{6}
Consider a function $y=f(x)=\sin x+\cos x+|\sin x-\cos x|$, $x\in[0,2\pi]$.
\qn{5} The value of $\displaystyle\lim_{h\to0}\frac{(f(\pi+h))^2-(f(\pi-h))^2}{h}$ is -
\qn{6} Number of points in $[0,2\pi]$ where $f(x)$ is non differentiable is -
\sect{SECTION-II(ii)}
\qn{7} Let $A=\begin{bmatrix}\frac12&\frac34\\1&-\frac12\end{bmatrix}$ then the value of $\operatorname{tr}(A)+\operatorname{tr}(A^2)+\operatorname{tr}(A^3)+\dots+\operatorname{tr}(A^{100})$, is\par
(Where $\operatorname{tr}(A)$ denotes trace of matrix $A$)
\qn{8} If the eccentricity of curve for which tangent at point $P$ intersects the $y$ axis at $M$ above $x$-axis such that origin is equidistant from $M$ and point of tangency is $e$ then value of $6e^2$ is ?
\qn{9} If $f(x)=e^{(x+1)^n}$; $(n\in N)$ and $f''(1)=67ke^k$ then value of $(n+k)$ is -
\raggedcolumns
\columnbreak
\papersec{2}
\partsec{1}{PHYSICS}
\sect{SECTION-I(i)}
\qn{1} A boy is jerking one end of a uniform taut string. The wave train propagating to the right has been shown in the figure.
\qfig{p2_phy_s1_q01_1}
\begin{qoptions}
\opt{A}{The boy is jerking at gradually decreasing frequency.}
\opt{B}{The boy is jerking at gradually increasing frequency.}
\opt{C}{Speed of particle at a is more than speed of particle at b.}
\opt{D}{Speed of particle at b is more than speed of particle at a.}
\end{qoptions}
\qn{2} A conducting loop has two straight and a circular arc sections. It is hinged at point $P$ and given $\omega_0=2$ rad/s angular velocity. Angle of circular arc is $\dfrac\pi4$ rad and radius 1m. Loop has total resistance $1\Omega$ and moment of inertia $\dfrac\pi4$ kgm$^2$. As loop enters in the magnetic field, induced current appears and hinge is insulator. Then :-
\qfig{p2_phy_s1_q02_1}
\begin{qoptions}
\opt{A}{Loop rotates uniformly for time $\dfrac\pi2$ s in magnetic field.}
\opt{B}{Loop rotates uniformly for time $\dfrac\pi4$ s in magnetic field.}
\opt{C}{Loop is able to complete more than 2 cycles if initial angular velocity is 8 rad/s.}
\opt{D}{Loop is able to complete more than 2 cycles if initial angular velocity is 7 rad/s.}
\end{qoptions}
\qn{3} A meterbridge circuit is shown in figure. The resistive wire has a total resistance of $6r$. Jockey divides the wire in two parts of length $a$ and $b$.
\qfig{p2_phy_s1_q03_1}
\begin{qoptions}
\opt{A}{For zero deflection in galvanometer $a:b=7:5$}
\opt{B}{For zero deflection in Galvanometer $a:b=4:3$.}
\opt{C}{If galvanometer is replaced by capacitor then for maximum time constant $a:b=7:5$}
\opt{D}{If galvanometer is replaced by capacitor then for maximum time constant $a:b=4:3$}
\end{qoptions}
\qn{4} A series R--L circuit and a series R--C circuit have same power factors $\dfrac1{\sqrt2}$. When these circuits are connected with same voltage source at different times, the maximum current of R-L circuit is found 9 times less than maximum current of R-C circuit. Now these circuits are connected in series together with same voltage source. Then :-
\begin{qoptions}
\opt{A}{Final circuit is capacitive in nature.}
\opt{B}{Final circuit is inductive in nature.}
\opt{C}{Power factor of final circuit is 0.78.}
\opt{D}{Power factor of final circuit is 0.62.}
\end{qoptions}
\qn{5} Three charge particles are placed at vertices of equilateral triangle. They start moving such that they always follow each other in anticlockwise manner along sides of triangle with uniform speed $v$. There is a superconducting ring at center of triangle having radius $r$. The nature of ring is such that its flux always remains constant with help of induced current. Ring is very small in size and charges were initially very far. Initially there is no current in ring.
\qfig{p2_phy_s1_q05_1}
\begin{qoptions}
\opt{A}{When side of triangle is $a$ the current in ring is $\dfrac{9qvr}{4\pi a^2}$.}
\opt{B}{When side of triangle is $a$ the current in ring is $\dfrac54\dfrac{qvr}{\pi a^2}$.}
\opt{C}{Magnetic dipole moment of ring varies linearly with time.}
\opt{D}{Magnetic force on any charge due to ring varies inversely with $5^{\text{th}}$ power of triangle's side.}
\end{qoptions}
\qn{6} A transparent sphere of refractive index 1.5 has radius $R$. There is a point light source $S$ inside the sphere at some distance from center. A person observes the light source along diametric line passing through $S$. Source appears at $\dfrac47R$ distance from surface.
\begin{qoptions}
\opt{A}{Source is at $\dfrac{2R}3$ distance from center of sphere.}
\opt{B}{When person observes the source from diametrically opposite point it appears $\dfrac{8R}5$ distance from surface.}
\opt{C}{All light from source passes the spherical boundary without any total internal reflection.}
\opt{D}{Source is at $\dfrac R3$ distance from center of sphere}
\end{qoptions}
\columnbreak
\sect{SECTION-I(ii)}
\paraq{7}{8}
Consider a fixed highly evacuated cylinder in a gravity free region. Inside the cylinder piston of mass $M$ can move without friction. In each part of the cylinder, a ball of mass $m$ ($m\ll M$) is moving back and forth with very high speed $v$ and in doing so they are colliding elastically with the ends of the cylinder and with the piston. When the piston is at the center, it is in equilibrium as shown in figure. Take $M=0.2$ kg, $m=10^{-3}$ kg $L=2$ m, $v=\pi$ m/s
\qfig{p2_phy_s2_para7_8_1}
\qn{7} Now the piston is shifted slightly towards one of the ends of cylinder with a negligible speed and released. Time period of oscillations of piston is :-
\begin{qoptions}
\opt{A}{1.5}
\opt{B}{20}
\opt{C}{3.2}
\opt{D}{45}
\end{qoptions}
\qn{8} If the ball on right side has double mass as well as speed of left side ball then length of right section of cylinder in steady state is :-
\begin{qoptions}
\opt{A}{$\dfrac{4L}5$}
\opt{B}{$\dfrac{8L}9$}
\opt{C}{$\dfrac{7L}8$}
\opt{D}{$\dfrac{3L}4$}
\end{qoptions}
\paraq{9}{10}
In resonance column experiment a tuning fork of frequency $f=400$ Hz is held above the pipe as shown in figure. The reservoir is raised and lowered to change the level of water and thus the length of the column of air in the tube. The area of cross section of the reservoir is 6 times that of the pipe. Initially, the reservoir is kept so that the pipe is full up to the brim. Tuning fork is sounded and the reservoir is lowred. When the reservoir is lowred by 21 cm, first resonance is recorded. When the reservoir is lowered further by 49 cm the second resonance is heard.
\qfig{p2_phy_s2_para9_10_1}
\qn{9} The area of reservoir is :-
\begin{qoptions}
\opt{A}{$54\pi$ cm$^2$}
\opt{B}{$112\pi$ cm$^2$}
\opt{C}{$27\pi$ cm$^2$}
\opt{D}{$150\pi$ cm$^2$}
\end{qoptions}
\qn{10} The speed of sound is :-
\begin{qoptions}
\opt{A}{228 m/s}
\opt{B}{336 m/s}
\opt{C}{398 m/s}
\opt{D}{328 m/s}
\end{qoptions}
\sect{SECTION-II(i)}
\paraq{1}{2}
A portion of a large network is shown in figure. Values of resistances are indicated and all ammeters are ideal. In steady state reading of ammeter $A_3$ is 3A.
\qfig{p2_phy_s3_para1_2_1}
\qn{1} Reading of ammeter $A_1$ is \ldots\ldots\ A.
\qn{2} Reading of ammeter $A_2$ is \ldots\ldots\ A.
\paraq{3}{4}
Large number of identical thin prisms are placed on $x$-axis. Each has refractive index 1.5 and angle of prism $\dfrac\pi2\times10^{-3}$rad. A light ray incidents parallel to $x$-axis along line $y=y_0$. Take thickness of each prism $t=1$ mm. Take $\ln2=0.7$.
\qfig{p2_phy_s3_para3_4_1}
\qn{3} When ray reaches $x=1$m angle made by ray with $x$-axis is $\theta$ rad. Value of $\theta$ is \ldots\ldots
\qn{4} When ray reaches $x=1$m its $y$ displacement magnitude in meter is \ldots\ldots\ldots
\paraq{5}{6}
An infinite conducting wire is fixed in gravity free space with a capacitor $C$. A conducting rod of mass $M$ and length $\ell$ can also slide on wire frame without any friction system is placed in uniform magnetic field $B$. Resistance of wire as well as rod is negligible and rod is initially at rest with respect to wire frame. Now rod is imparted an impulse towards positive $x$ axis $\vec J=J\hat i$ at $t=0$. Take $J=4$ kg m/s, $M=1$kg, $\ell=0.5$ m, $B=2$T, $C=2$F.
\qfig{p2_phy_s3_para5_6_1}
\qn{5} Displacement of rod in $t=1$s is \ldots\ldots\ldots
\qn{6} Energy lost in circuit till steady state is achieved is \ldots\ldots
\sect{SECTION-II(ii)}
\qn{7} In the circuit shown in figure. Switch is initially open for long time and at $t=0$ it is closed. Current in switch at $t=16$ s is $\dfrac N3$. Find value of $N$. Ignore the mutual inductance.
\qfig{p2_phy_s4_q07_1}
\qn{8} A travelling wave is propagating along a tight rope. Velocity of any point of rope is shown in the graph. There is a point $P$ on rope. If frequency of wave is $\dfrac{\sqrt3}4$Hz then magnitude of acceleration of $P$ is ?
\qfig{p2_phy_s4_q08_1}
\qn{9} A string of mass 1gm, length 3m, tension 30 N is fixed at both the ends and is oscillating in 3rd overtone. The amplitude at the antinode is 0.1 mm. What is the energy of oscillations (in $\mu$J) ?
\partsec{2}{CHEMISTRY}
\sect{SECTION-I(i)}
\qn{1} Which of the following statement(s) is/are incorrect?
\begin{qoptions}
\opt{A}{Nickel salts gives green borax bead test in oxidising flame.}
\opt{B}{From mixture of AgCl and AgI, NH$_4$OH solution dissolves only AgCl}
\opt{C}{Fe$^{3+}$ ions given a deep green precipitate, on adding K$_4$[Fe(CN)$_6$] solution}
\opt{D}{On strong heating Na$_2$S$_2$O$_3$.5H$_2$O, a polysulphide is formed}
\end{qoptions}
\qn{2} During extraction of Fe, Select the incorrect statement(s).
\begin{qoptions}
\opt{A}{In reduction zone, Fe$_2$O$_3$ is mainly reduced by carbon}
\opt{B}{CO$_2$ is one of the by-products of reduction zone of blast furnace.}
\opt{C}{Cast iron is obtained from blast furnace.}
\opt{D}{Uppermost zone of blast furnace is slag formation zone.}
\end{qoptions}
\qn{3} Consider the cell Pb(s) $|$ PbI$_2$(s) $|$ I$^-$(aq. 0.1M) $\|$ Cu$^{2+}$(aq. 0.01M) $|$ Cu(s) at temperature $T$.\par
Given $K_{\text{sp PbI}_2}=10^{-8}$ and $E^\circ_{\text{Pb}^{2+}/\text{Pb}}=-0.13$V, $E^\circ_{\text{Cu}^{2+}/\text{Cu}}=0.34$V and $\dfrac{2.303\,RT}{F}=0.06$V.\par
Select the correct option(s) :
\begin{qoptions}
\opt{A}{Above cell is nonspontaneous}
\opt{B}{$E_{\text{cell}}=-0.59$ V}
\opt{C}{$[\text{Pb}^{2+}]_{\text{Anode}}=10^{-6}$ M}
\opt{D}{$E_{\text{cell}}=0.59$ V.}
\end{qoptions}
\qn{4} Compressibility of real gas will be less than ideal gas :
\begin{qoptions}
\opt{A}{At low pressure when, temperature of gas $<$ Boyle's temperature}
\opt{B}{At low pressure when, temperature of gas $>$ Boyle's temperature}
\opt{C}{At very high pressure when, temperature of gas $<$ Boyle's temperature}
\opt{D}{At very high pressure when, temperature of gas $>$ Boyle's temperature}
\end{qoptions}
\qn{5} Which of the following aicd remains affected on heating?
\begin{qoptions}
\opt{A}{Fumaric acid}
\opt{B}{Malonic acid}
\opt{C}{Maleic acid}
\opt{D}{Succinic acid}
\end{qoptions}
\qn{6}
\qfig{p2_chem_s1_q06_1}
Select correct statement :
\begin{qoptions}
\opt{A}{Q is \qfigi{p2_chem_s1_q06_2}}
\opt{B}{R is \qfigi{p2_chem_s1_q06_3}}
\opt{C}{S is \qfigi{p2_chem_s1_q06_4}}
\opt{D}{Q is \qfigi{p2_chem_s1_q06_5}}
\end{qoptions}
\sect{SECTION-I(ii)}
\paraq{7}{8}
Two cations on reacting with group-2 reagent gives two different colour ppt. `X' and `Y'.\par
\qn{7} X + conc. HCl $\xrightarrow{\Delta}$ W + gas 1\par
W + H$_2$O $\rightleftharpoons$ R (white ppt.)\par
R is not soluble in tartaric acid, colour of precipitate `X' is
\begin{qoptions}
\opt{A}{Black}
\opt{B}{Yellow}
\opt{C}{Orange}
\opt{D}{Green}
\end{qoptions}
\qn{8} Y + conc. HCl $\xrightarrow{\Delta}$ Z + Gas 1\par
Z + H$_2$O $\rightleftharpoons$ Q (white ppt.)\par
Q is soluble in tartatic acid, colour of precipitate `Y' is :-
\begin{qoptions}
\opt{A}{Black}
\opt{B}{Yellow}
\opt{C}{Orange}
\opt{D}{Green}
\end{qoptions}
\paraq{9}{10}
\qfig{p2_chem_s2_q08_1}
\qn{9} Which of the following reaction does not take place in fusion zone (of hearth)
\begin{qoptions}
\opt{A}{MnO + C $\to$ Mn + CO$\uparrow$}
\opt{B}{SiO$_2$ + 2C $\to$ Si + 2CO}
\opt{C}{2P$_2$O$_5$ + 10C $\to$ 4P + 10CO$\uparrow$}
\opt{D}{Fe$_2$O$_3$ + 3CO $\to$ 2Fe + 3CO$_2$}
\end{qoptions}
\qn{10} Choose the correct option regarding `x'.\par
Statement-1 : In the process `x' impurities of Mn and Si are mainly removed in the form of MnSiO$_3$\par
Statement-2 : In process `x', Ca$_3$(PO$_4$)$_2$ may be produced, which is a fertilizer.\par
Statement-3 : In molten `Fe' affinity of O$_2$ is in order Si $>$ Mn $>$ P $>$ Fe $>$ C\par
Statement-4 : Spiegeleisen is an alloy of Fe, C and Mg which is used in steel making.
\begin{qoptions}
\opt{A}{TTTF}
\opt{B}{TFTT}
\opt{C}{FFTT}
\opt{D}{TTFF}
\end{qoptions}
\sect{SECTION-II(i)}
\paraq{1}{2}
Consider the combustion of 1 mole of CH$_4$ according to the following reaction with 20\% excess air at $27^\circ$C and 1 bar pressure.\par
\noindent\hfill CH$_4$(g) + 2O$_2$(g) $\to$ CO$_2$(g) + 2H$_2$O(g)\hfill\null\par
Assume air as a mixture of 20\% O$_2$ and 80\% N$_2$ by volume.\par
Given : $\Delta H^\circ_f$, CO$_2$(g) $=-100$ K cal mol$^{-1}$\par
$\Delta H^\circ_f$, H$_2$O(g) $=-60$ K cal mol$^{-1}$\par
$\Delta H^\circ_f$, CH$_4$(g) $=-20$ K cal mol$^{-1}$\par
Assume all the gases are ideal gases and vibrational degree of freedom is not active. $R=2$ cal K$^{-1}$ mol$^{-1}$.
\qn{1} Calculate maximum temperature in Kelvin that can be reached by above combustion. Ignore the variation of $\Delta_rH$ with temperature.
\qn{2} If energy released by combustion of 1 mole of CH$_4$ at $127^\circ$C and 1 bar is supplied to a heat engine at $127^\circ$C which discards heat at $27^\circ$C and whose efficiency is 80\% of maximum possible efficiency then calculate amount of heat discarded by the heat engine in Kcal. ($R=2$ cal K$^{-1}$ mol$^{-1}$).
\paraq{3}{4}
A weak base BOH is titrated with a strong acid HA. When 10 ml of HA is added, the pH is 10.3 and when 25 ml of HA is added the pH is 9.3.\par
In another experiment 40g of the salt BA is dissolved in water to make 1L aq. solution. The pH of solution is found to be 5.15. If salt BA forms NaCl type crystal having $r_{\text{B}^+}=0.8$\AA\ and $r_{\text{A}^-}=1.2$\AA\ ($\log2=0.3$, $N_A=6\times10^{23}$).
\qn{3} Find p$K_b$ of BOH :
\qn{4} Find density of BA crystal in g/ml.
\paraq{5}{6}
\qfig{p2_chem_s3_para5_6_1}
\qn{5} X react with hot and concentrated Y to give chloride and Z\par
Find total resonating structure of anion of Z.
\qn{6} How many carbon present in principle carbon chain of product (D) according to IUPAC rule?
\sect{SECTION-II(ii)}
\qn{7} How many of the following will have yellow ppt. as the product.\par
(i) Pb$^{2+}$ (aq.) + 2NaOH $\to$\par
(ii) PbCrO$_4$ + dil. HNO$_3\to$\par
(iii) Ag$_2$CrO$_4$ + excess NH$_3\to$\par
(iv) Hg$^{2+}$ (aq.) + NaOH $\to$\par
(v) K$_2$[Cd(CN)$_4$] + H$_2$S $\to$\par
(vi) Ni$^{2+}$ (aq.) + CN$^-$ (excess) $\to$\par
(vii) H$_2$SO$_4$ (aq.) + (NH$_4$)$_2$CrO$_4\to$\par
(viii) Pb$^{+2}$ (aq.) + (NH$_4$)$_2$CrO$_4\to$
\qn{8} How many isomers ``X'' of C$_8$H$_{10}$ when reacts with hot alkaline KMnO$_4$ give only aromatic dicarboxylic acid?
\qn{9} How many Isomers Y of C$_4$H$_8$ when reacts with hot alkaline KMnO$_4$ give carbondioxide?
\partsec{3}{MATHEMATICS}
\sect{SECTION-I(i)}
\qn{1} Consider a square matrix $A$. Identify the correct statement -
\begin{qoptions}
\opt{A}{If $A$ is nilpotent matrix of index 2 then $A(I+A)^n=A\ \forall\ n\in N$}
\opt{B}{If $A$ is nilpotent matrix of index 2 then $A(I-A)^n=A\ \forall\ n\in N$}
\opt{C}{If $A$ is idempotent matrix, then $(I+A)^n=I+(2^n-1)A\ \forall\ n\in N$}
\opt{D}{If $A$ is idempotent matrix, then $(I+A)^n=I-(2^n-1)A\ \forall\ n\in N$}
\end{qoptions}
\qn{2} If $z=\log\tan\dfrac t2$, and $\dfrac{d^2y}{dz^2}+4y=0$, then find the correct option
\begin{qoptions}
\opt{A}{$\dfrac{d^2z}{dt^2}+\operatorname{cosec}t\cot t=0$}
\opt{B}{$\dfrac{d^2y}{dt^2}+\cot t\dfrac{dy}{dt}+4y\operatorname{cosec}^2t=0$}
\opt{C}{$\dfrac{d^2y}{dt^2}+\cot t\dfrac{dy}{dt}+2y\operatorname{cosec}t=0$}
\opt{D}{$\dfrac{d^3y}{dt^3}+\cos t\dfrac{dy}{dt}+y\cot t=0$}
\end{qoptions}
\qn{3} Let $f(x)=\displaystyle\int_x^{x+\pi/3}|\sin\theta|\,d\theta$; where $x\in[0,\pi]$, which of the following is/are true
\begin{qoptions}
\opt{A}{maximum value of $f(x)$ is 1}
\opt{B}{minimum value of $f(x)$ is $\dfrac32$}
\opt{C}{maximum value of $f(x)$ is $\dfrac12$}
\opt{D}{minimum value of $f(x)$ is $2-\sqrt3$}
\end{qoptions}
\qn{4} Let $f$ is a differentiable function satisfying $3f(x+2y)=f(2x)+f(y)+2x+10y+1\ \forall\ x,y\in R$ such that $f'(0)=2f(0)$, then - (where $C$ is constant of integration)
\begin{qoptions}
\opt{A}{$f(x)$ is an odd function}
\opt{B}{$\displaystyle\int\frac{dx}{f(x)}=\ell n\left(\sqrt{|f(x)|}\right)+C$}
\opt{C}{$f(0.5)=f'(5.0)$}
\opt{D}{$\sin(f(x))$ is a periodic function with period $\pi$}
\end{qoptions}
\qn{5} Choose the correct statement(s) -
\begin{qoptions}
\opt{A}{The largest integer in the range of function $(\sin^{-1}x)^2-(\cos^{-1}x)^2$ is 2}
\opt{B}{The smallest integer in the range of function $(\sin^{-1}x)^2-(\cos^{-1}x)^2$ is $-7$}
\opt{C}{The value of `$x$' satisfying $\sin^{-1}x+2\cos^{-1}(-x)=2\pi$, is greater than $\dfrac12$}
\opt{D}{The value of `$x$' satisfying $\sin^{-1}x+2\cos^{-1}(-x)=2\pi$, is less than $\dfrac12$}
\end{qoptions}
\qn{6} Let $f_1:(-\infty,0)\to R$, $f_2:(-\infty,0)\to R$ and $f_3:(1,\infty)\to R$ be functions defined by\par
$f_1(x)=e^x-\sqrt{-x}$, $f_2(x)=\sqrt{-x}+|x|$ and $f_3(x)=e^{1-x}-\sqrt{x-1}$. Then choose the correct option(s)
\begin{qoptions}
\opt{A}{$f_1(x)$ is increasing function but $f_2(x)$ and $f_3(x)$ are decreasing functions}
\opt{B}{$f_1(x)$ is decreasing function but $f_2(x)$ is increasing function}
\opt{C}{$f_2(x)$ and $f_3(x)$ are decreasing functions}
\opt{D}{$f_2(x)$ is decreasing function but $f_1(x)$ and $f_3(x)$ are increasing functions}
\end{qoptions}
\sect{SECTION-I(ii)}
\paraq{7}{8}
Let $f(x)$ be a differentiate function such that $f'(x)=\sqrt{2f(x)-f^2(x)}$, where $f'(0)=f(0)=1$
\qn{7} Area bounded by $y=f(x)$ and $x$-axis from $x=0$ to $x=\pi$ is
\begin{qoptions}
\opt{A}{$2\pi+4$}
\opt{B}{$\pi+2$}
\opt{C}{$\pi+4$}
\opt{D}{$2\pi+2$}
\end{qoptions}
\qn{8} $\displaystyle\int_0^{\pi/2}\frac{f(x)-1}{\cos^2x+1}\,dx$ is
\begin{qoptions}
\opt{A}{$\dfrac\pi4$}
\opt{B}{$-\dfrac\pi4$}
\opt{C}{$\dfrac\pi2$}
\opt{D}{$\dfrac{3\pi}4$}
\end{qoptions}
\paraq{9}{10}
Consider $f(x)=x^{\ell nx}$ and $g(x)=e^2x$. Let $\alpha$ and $\beta$ ($\alpha<\beta$) be two values of $x$ satisfying equation $f(x)=g(x)$.
\qn{9} Product of $\alpha\beta$ is -
\begin{qoptions}
\opt{A}{$\dfrac1{e^2}$}
\opt{B}{$\dfrac1e$}
\opt{C}{$e$}
\opt{D}{$e^2$}
\end{qoptions}
\qn{10} If $\displaystyle\lim_{x\to\beta}\frac{f(x)-c\beta}{g(x)-\beta^2}$ exists and equal to `$r$' then value of $(c-r)$ is -
\begin{qoptions}
\opt{A}{$4-e^2$}
\opt{B}{$e^2-4$}
\opt{C}{$4-e$}
\opt{D}{$e-4$}
\end{qoptions}
\columnbreak
\sect{SECTION-II(i)}
\paraq{1}{2}
Let $A=\begin{bmatrix}1&0&0\\1&0&1\\0&1&0\end{bmatrix}$ satisfy $A^n=A^{n-2}+A^2-I\ \forall\ n\ge3$. Further consider a matrix $U$ of order $3\times3$ with its column as $U_1,U_2,U_3$ such that $A^{50}U_1=\begin{bmatrix}1\\25\\25\end{bmatrix}$; $A^{50}U_2=\begin{bmatrix}0\\1\\0\end{bmatrix}$; $A^{50}U_3=\begin{bmatrix}0\\0\\1\end{bmatrix}$
\qn{1} Trace of $A^{50}$ is
\qn{2} Sum of all entries of $UA^{50}U^{-1}$ is
\paraq{3}{4}
Let $f(x)=4^x$ and $g(x)=\sin(\cos(\cos\pi x)+\cos(\sin\pi x))$, then
\qn{3} The value of $\displaystyle\lim_{x\to\infty}\frac{g(x)}{f(x)}$ is equal to
\qn{4} $\displaystyle\int_0^4f(x)\cdot g(x)\,dx=\lambda\int_0^{1/2}f(x)\,g(x)\,dx$, then $\lambda$ is -
\paraq{5}{6}
Let $y=f(x)$ be a fourth degree algebraic polynomial function such that\par
(i) $\displaystyle\lim_{x\to1}\frac{f(x)}{(x-1)^2}=2$\par
(ii) It cuts $y$-axis a point where slope of normal is $\dfrac18$\par
(iii) Tangent at $x=2$ is parallel to $y=8x+1$.\par
Consider $y=g(x)$ be a quadratic curve touching and cutting $y=f(x)$ on $x$-axis and $y$-axis respectively.
\qn{5} If angle of intersection between $y=f(x)$ \& $y=g(x)$ at point where they cut $y$-axis is $\tan^{-1}\dfrac ab$ (where $a$ and $b$ are co-prime) then value of $(a+b)$ is -
\qn{6} The abscissae of point on curve $y=g(x)$, where length of subtangent is 2018 can be -
\sect{SECTION-II(ii)}
\qn{7} If area bounded on the left by $y$-axis, below by $x$-axis, right by $x=\dfrac\pi2$ above left by $y=\cos x$ and above right by $y=\sin x$ is $K$ then value of $100K^2$ is -
\qn{8} Let $A$ and $B$ be two sets such that $n(A)=4$ and $n(B)=3$. If $(x,1)$, $(y,2)$, $(z,1)$ and $(2,3)$ are in $A\times B$, where $x$, $y$, $z$ are distinct elements. If a function $f:A\to B$ defined \& number of onto function from set $A$ to set $B$ is $\lambda$, then value of $\dfrac\lambda6$ is equal to
\qn{9} Let matrix $P=\begin{bmatrix}\frac1{\sqrt2}&\frac1{\sqrt2}\\-\frac1{\sqrt2}&\frac1{\sqrt2}\end{bmatrix}$ and $Q$ is another matrix such that $Q^TP=P^T$ and $\lambda Q^T=2P^T-\sqrt2I$, then value of determinant of matrix $\operatorname{adj}(I+Q^T)\cdot\operatorname{adj}(\lambda P)$ is\par
(where $I$ is unit matrix of order 2, $\operatorname{adj}(X)$ and $X^T$ denote adjoint and transpose of matrix $X$ respectively)
\end{multicols}
\end{document}
